% Recap and reference quantities --- plan \S4.2 (Phase 1).
% Sources: new_notes3 Part A (\S\S4--7), new_notes Part I \S1--\S5.
% All boxed formulas are the CORRECTED forms (see errata remarks).

\section{The single-cell BDC process: recap and reference quantities}

Chapter~3 defined the birth--death--catastrophe process and derived its
principal quantities. This chapter builds on that foundation, and several of
the results it needs are restated here --- in corrected and simplified form ---
so that what follows is self-contained. Section~\ref{sec:notation-tab} fixes
the notation used throughout; the remaining sections collect the closed forms
that will be quoted repeatedly. Derivations are sketched; the full accounts
live in Chapter~3 and, for the corrected items, in the remarks.

\subsection{Canonical notation}
\label{sec:notation-tab}

The notation of this chapter is collected in Table~\ref{tab:notation}. It
resolves the collisions that arise the moment the intracellular and
population models are joined: $\beta$ vs $\lambda$ for the birth rate;
$\delta$ for the catastrophe rate vs the infected-cell removal rate; $I$ and
$V$ for probabilities and expectations vs cell and virion counts; $H$ vs $R$
for the rupture state.

\begin{table}[H]
\centering
\renewcommand{\arraystretch}{1.25}
\small
\begin{tabular}{@{}llp{0.44\textwidth}@{}}
\toprule
Symbol & Meaning & Notes \\
\midrule
$\beta,\mu,\delta$ & birth, death, catastrophe rates (per capita) & $\lambda$ is not used for birth \\
$\xt$ & intracellular count (CTMC) & \\
$R$ & absorbing rupture state & \\
$\wt$ & released count from one cell & $0$ until burst \\
$\tau$ & burst time (possibly $\infty$) & \\
$\Kb$ & burst size (random variable) & never bare $K$, since $K(t)=\mean{\xt^2}$ \\
$I(t)$ & $\Prob{\text{no burst by }t}$ & not an infected-cell count \\
$D(t)$ & $\Prob{\text{cleared internally, no burst}}$ & $= p_0(t)$ \\
$\Ihat(t)$ & $\Prob{\xt\notin\{0,R\}} = I(t)-D(t)$ & productive infection \\
$J(t),K(t)$ & $\mean{\xt}$, $\mean{\xt^2}$ & \\
$V(t)$ & $\mean{\wt}$, cumulative mean release & not a free-virion count \\
$a,b$ & characteristic roots, $b<1<a$ & \\
$A,B,\theta$ & $a-1$, $1-b$, $\beta(a-b)$ & \\
$L$ & $a(1-b)=a-\mu/\beta$ & flooding parameter \\
$T$ & target-cell count (held fixed) & population model \\
$\gamma$ & infection rate & not $\beta$ \\
$\Icell,\Vfree$ & infected cells, free virions & population counts \\
$c$ & virion clearance & \\
$d_{\Icell}$ & infected-cell removal (phenomenological) & not $\delta$ \\
$p$, $p_{\rm eff}(r)$ & constant / effective release rate & \\
$g(a)$ & release kernel (expected flux at age $a$) & \\
$q$ & $\gamma T/(\gamma T+c)$, infection success & branching \\
\bottomrule
\end{tabular}
\caption{Canonical notation for this chapter.}
\label{tab:notation}
\end{table}

\subsection{Process definition and the joint process}

One cell, one intracellular population. The count $\xt$ is a continuous-time
Markov chain on $\{0,1,2,\ldots\}\cup\{R\}$ which, from any state $n\ge1$,
jumps
\begin{align*}
n &\to n+1 &&\text{at rate } \beta n \quad\text{(birth)},\\
n &\to n-1 &&\text{at rate } \mu n \quad\text{(death)},\\
n &\to R &&\text{at rate } \delta n \quad\text{(catastrophe)}.
\end{align*}
Both $0$ and $R$ are absorbing: the population either clears itself
internally, or the cell ruptures. Unless stated otherwise we start from a
single founder, $X_0=1$. The process of classical references
\cite{brockwell1982birth,karlin1982linear,di2008note} is slightly more
general (immigration, distributed catastrophe sizes); the specialisation
above is the one with the closed-form family we need.

The rupture event is what matters for applications, so we record the release
count alongside the load. Let $\tau$ be the time of catastrophe (possibly
$\tau=\infty$ if the population clears internally first). The released count
$\wt$ is
\[
\wt = \begin{cases}
0, & t<\tau,\\
X_{\tau^-}, & t\ge\tau,
\end{cases}
\]
that is, at rupture the whole pre-burst load is dumped and then held. All
single-cell release quantities in this chapter are functionals of the joint
process $(\xt,\wt)$.

\subsection{Roots and the workhorse identity}

Write
\begin{equation}
\eta=\frac{\beta+\mu+\delta}{2\beta},\qquad
a=\eta+\sqrt{\eta^2-\tfrac{\mu}{\beta}},\qquad
b=\eta-\sqrt{\eta^2-\tfrac{\mu}{\beta}},
\end{equation}
the two roots of the quadratic $\beta f^2-(\beta+\mu+\delta)f+\mu$ that will
appear in the next subsection, and set
\begin{equation}
A:=a-1>0,\qquad B:=1-b>0,\qquad \theta:=\beta(a-b)
=\sqrt{(\beta+\mu+\delta)^2-4\beta\mu}.
\end{equation}
Vieta's relations give $ab=\mu/\beta$ and $a+b=(\beta+\mu+\delta)/\beta$, and
the identity that will do most of the algebraic work in this chapter is
\begin{equation}
AB=(a-1)(1-b)=(a+b)-ab-1=\frac{\delta}{\beta},
\qquad A+B=a-b.
\label{eq:AB}
\end{equation}
That $b<1<a$ always is immediate from
$\beta(1-a)(1-b)=\beta\lrb{1-(a+b)+ab}=-\delta<0$, which places $1$ strictly
between the two roots.

\subsection{Three fixation functions}

There are three absorbing-behaviour probabilities, not two, and it pays to
name all of them:
\begin{definition}
\begin{equation}
I(t)=\Prob{\text{no burst by }t},\qquad
D(t)=\Prob{\text{internal extinction by }t,\ \text{no burst}},
\end{equation}
\begin{equation}
\Ihat(t)=\Prob{\xt\notin\{0,R\}}=I(t)-D(t),
\end{equation}
the last being the probability that the cell is still productively infected.
\end{definition}

Both $I$ and $D$ satisfy the same Riccati equation. For $I$, condition on the
first event of the founder: a birth (rate $\beta$) replaces the one lineage
by two independent ones, both of which must avoid bursting, contributing
$I^2$; a death (rate $\mu$) makes bursting impossible, contributing $1$; a
catastrophe (rate $\delta$) contributes $0$. Hence
\begin{equation}
\ddt{I}=\beta\lrb{I^2-I}+\mu\lrb{1-I}+\delta\lrb{0-I}
=\beta(I-a)(I-b)=\mu-(\beta+\mu+\delta)I+\beta I^2,
\label{eq:riccati}
\end{equation}
with $I(0)=1$. The same first-step decomposition with ``cleared internally,
no burst'' in place of ``no burst'' gives the identical equation for $D$,
now with $D(0)=0$. Solving \eqref{eq:riccati} by separation and writing
$w:=\ee^{\theta t}\ge1$,
\begin{equation}
I(t)=\frac{aB+bAw}{B+Aw},\qquad
D(t)=\frac{ab\,(w-1)}{aw-b},\qquad
\Ihat(t)=\frac{(a-b)^2\,w}{(B+Aw)(aw-b)}.
\label{eq:IDIhat}
\end{equation}
The limits are $I(\infty)=D(\infty)=b$, $\Ihat(\infty)=0$, with
$\Ihat(0)=1$ and $\Ihat'(0)=-(\mu+\delta)$ --- at the instant of founding,
the productive state is lost only by death or catastrophe of the single
founder.

\begin{remark}[Correction of the Chapter~3 formula for $\Ihat$]
\label{rem:ihat}
Chapter~3 gives
$\Ihat(t)=\hat\eta/\lrb{1-(1-\hat\eta)\ee^{\hat\eta\beta t}}$ with
$\hat\eta=(\beta+\mu+\delta)/\beta$, derived from the ODE
$\Ihat'=-(\beta+\mu+\delta)\Ihat+\beta\Ihat^2$. The branching term is wrong:
on a birth, the two daughter lineages must avoid bursting \emph{and} not both
go extinct, an event of probability $I^2-D^2$, not $\Ihat^2$. The error is
invisible when $\mu=0$ (then $D\equiv0$ and $\Ihat=I$), which is presumably
why it survived. Numerically, at $(\beta,\mu,\delta)=(1,0.2,0.05)$ and
$t=1$: the correct value is $\Ihat\approx0.8087$; the thesis formula gives
$\approx0.6675$. An independent confirmation arrives in
\S\ref{sec:state-probs}: the state probabilities $p_n(t)$ sum over $n\ge1$
to exactly the right-hand side of \eqref{eq:IDIhat}.
\end{remark}

\subsection{Moments and mean release}
\label{sec:moments-recap}

The hazard of bursting at time $t$ is $\delta$ times the current load, so
$I'(t)=-\delta J(t)$; inserting the Riccati equation \eqref{eq:riccati}
then gives the mean
\begin{equation}
J(t)=\mean{\xt}=-\frac{1}{\delta}\ddt{I}
=\frac{(a-b)^2\,w}{(B+Aw)^2}
=-\frac{\beta}{\delta}(I-a)(I-b),
\label{eq:J}
\end{equation}
with $J(0)=1$ now manifest. The second moment,
\begin{equation}
K(t)=\mean{\xt^2}=\Bigl[1+\tfrac{2\beta}{\delta}(1-I)\Bigr]J
=\frac{2\beta}{\delta}(\kappa-I)J,\qquad \kappa=1+\frac{\delta}{2\beta},
\label{eq:K}
\end{equation}
follows from the moment hierarchy: one checks by differentiation, using
\eqref{eq:riccati}, that the right-hand side satisfies
$K'=2(\beta-\mu)K+(\beta+\mu)J-\delta\mean{\xt^3}$ with the third moment
eliminated exactly as in Chapter~3.

The cumulative mean release $V(t)=\mean{\wt}$ increases only at rupture
events, and a rupture from state $n$ occurs at rate $\delta n$ and releases
$n$ particles, so the expected release flux is $\delta\mean{\xt^2}$:
\begin{equation}
V'(t)=\delta K(t),\qquad V(0)=0.
\end{equation}
Integrating with \eqref{eq:J}--\eqref{eq:K} yields the simplification
\begin{equation}
V(t)=(1-I)\Bigl[1+\tfrac{\beta}{\delta}(1-I)\Bigr]
=\frac{\beta}{\delta}\lrb{(I-\kappa)^2-(1-\kappa)^2},
\label{eq:V}
\end{equation}
a quadratic in $(1-I)$ alone --- no $J$, no separate $\mu$. Chapter~3's form
$V=\frac{\beta-\mu}{\delta}(1-I)-J+1$ is the same quantity, less happily
dressed; the prefactor that clings to $J$ in Chapter~3 cancels identically by
\eqref{eq:AB}.

The lifetime mean release and the conditional mean burst size are
\begin{equation}
V_\infty=\frac{a(1-b)}{a-1}=\frac{aB}{A}
=\frac{\beta-\mu}{\delta}(1-b)+1,
\qquad
\mean{\Kb\mid\text{burst}}=\frac{V_\infty}{1-b}=\frac{a}{a-1}.
\label{eq:Vinf}
\end{equation}
A small point with large consequences: $V_\infty$ is
$\mean{\Kb\mathbf 1\{\text{burst}\}}$, and so \emph{includes} the
realisations that cleared internally and released nothing (a mass $b$ at
$0$). The conditional mean removes them. When $\mu=0$: $b=0$,
$a=1+\delta/\beta$, and $V_\infty=1+\beta/\delta$ coincides with the
conditional mean --- a coincidence that, as we shall see, one should not
lean on.

\subsection{Second moment of the release}
\label{sec:EW2-sec}

For the population-level variance in \S\ref{sec:discussion-assay} we will
need $\mean{\wt^2}$. The release variable jumps once, from $0$ to
$X_{\tau^-}$, so $\ddt{}\mean{\wt^2}=\delta\mean{\xt^3}$. The third moment
is eliminated via the second-moment ODE,
$K'=2(\beta-\mu)K+(\beta+\mu)J-\delta\mean{\xt^3}$, giving
\[
\ddt{}\mean{\wt^2}=2(\beta-\mu)K+(\beta+\mu)J-K'
=\ddt{}\lrb{\frac{2(\beta-\mu)}{\delta}V-\frac{\beta+\mu}{\delta}I-K},
\]
where the second equality uses $V'=\delta K$ and $J=-\delta^{-1}I'$ --- note
the sign. Integrating from $0$ with
$(I,J,K,V,\mean{W^2})(0)=(1,1,\mathbf 1,0,0)$ --- in particular $K(0)=\mean{X_0^2}=1$,
not $0$ ---
\begin{equation}
\mean{\wt^2}
=\frac{2(\beta-\mu)}{\delta}V
-\frac{\beta+\mu}{\delta}I
-K
+\frac{\beta+\mu}{\delta}
+1,
\label{eq:EW2}
\end{equation}
and $\mathrm{Var}(\wt)=\mean{\wt^2}-V^2$.

\begin{remark}[Correction and check]
\label{rem:EW2}
Chapter~3 \S8 has wrong signs on $I$ and $K$ and omits the constant $+1$;
both errors trace to the two slips above (the sign of $J=-\delta^{-1}I'$ and
the value of $K(0)$). An independent verification: take $\mu=0$ and
$t\to\infty$. Then $I,K\to0$ and $V\to1+\beta/\delta$, so \eqref{eq:EW2}
gives $\mean{W_\infty^2}=2\beta^2/\delta^2+3\beta/\delta+1$. But when
$\mu=0$ the burst size is geometric on $\{1,2,\ldots\}$ with success
probability $s=\delta/(\beta+\delta)$ (proved in \S\ref{sec:burst-size}),
whose second moment is $(2-s)/s^2=2\beta^2/\delta^2+3\beta/\delta+1$. Exact
agreement.
\end{remark}
